\section{A codimension bound over the coefficient ring}
\label{sec:codimension-stabilization}
The degree needed to generate a finite algebra can be bounded by the
codimension of the initial subspace, rather than by the length of the
algebra. A fibrewise dimension argument suggests this bound but does
not prove the assertion over a nonreduced parameter scheme. The
following argument proves an equality of submodules over the
coefficient ring.

\begin{theorem}[Codimension stabilization]
\label{thm:codimension-stabilization}
Let $B$ be a commutative finite locally free algebra over a
$\mathbb Q$-scheme, and let $W\subset B$ be a subbundle containing
$1$, with $B/W$ locally free of rank $r$. Write $W^j$ for the image
of $\Sym^j W\to B$. Then
\begin{equation}\label{eq:codimension-stabilization}
 W^{r+1}=\mathcal O[W],\qquad W^j=W^{r+1}\quad(j\ge r+1).
\end{equation}
These equalities remain true after every base change, without a
flatness hypothesis on $W^j$.

For a generating rank-$k$ subbundle $A$ of a rank-$d$ algebra over
a $\C$-scheme, all Fitting ideals of the multiplication cokernel
are consequently independent of the degree $j\ge d-k+1$.
The generating opens and normalization by units are those of
Lemma~\ref{lem:generating-open}. This is a statement about ideal
equalities; the unnormalized cokernels need not have a preferred
isomorphism in different degrees.
\end{theorem}
\begin{proof}
The case $r=0$ is immediate. Work otherwise on a free chart and
choose a splitting $B=W\oplus M$, where $\rank M=r$. For $w\in W$
write multiplication by $w$ in blocks
\[
 L_w=\begin{pmatrix}A_w&D_w\\ \beta_w&C_w\end{pmatrix}.
\]
Thus $\beta_w:W\to M$ and $C_w:M\to M$. Let
$N_0=\sum_{w\in W}\operatorname{im}\beta_w$, and, for $j\ge0$,
let $N_j$ be the sum of the images of $N_0$ under all words of length
at most $j$ in the operators $C_w$. A finite local basis of $W$
suffices in these definitions. Since $1\in W$, every $W^m$ contains
$W$. Projection to $M$ and subtraction of its $W$ component therefore
give, as submodules of the chosen splitting,
\begin{equation}\label{eq:residual-filtration}
 W^{j+2}=W\oplus N_j.
\end{equation}
Indeed the initial case is the definition of $N_0$, and multiplication
by $W$ changes the residual module $N$ into
$N_0+\sum_w C_w N$.

Commutativity of the full multiplication operators gives the identity
\begin{equation}\label{eq:compressed-commutator}
 [C_w,C_v]=\beta_vD_w-\beta_wD_v,
 \qquad \operatorname{im}[C_w,C_v]\subset N_0.
\end{equation}
Consider a word of length $r$ applied to $N_0$. Interchanging two
adjacent factors changes it by a word with a commutator in those
positions. The commutator sends its entire argument into $N_0$;
at most $r-2$ factors remain to its left. The change thus belongs
to $N_{r-2}$, and in particular to $N_{r-1}$. For $r=1$ no
interchange is needed. All permutations of a length-$r$ word
therefore agree modulo $N_{r-1}$ when applied to $N_0$.

Apply Cayley--Hamilton to the $r$ by $r$ matrix
$C(t)=\sum_{i=1}^r t_i C_{w_i}$ and extract the coefficient of
$t_1\cdots t_r$. The leading term is the sum of the $r!$ ordered
products of the $C_{w_i}$. Every other term is a scalar linear
combination of words of length at most $r-1$. Applied to $N_0$
these other terms lie in $N_{r-1}$. The preceding interchange
argument and the invertibility of $r!$ now show that every
length-$r$ word sends $N_0$ into $N_{r-1}$. Hence
$N_r=N_{r-1}$. This equality says that $N_{r-1}$ is invariant under
all $C_w$, so the residual filtration is constant thereafter.
Equation~\eqref{eq:residual-filtration} proves
\eqref{eq:codimension-stabilization}.

This calculation used no division except by $r!$ and no passage
to residue fields. It is valid over every base-changed coefficient
ring; in particular, it does not assert that an unsplit image is a
vector bundle. Finally choose a unit $a\in A$ smoothly locally
and put $W=a^{-1}A$. Multiplication by $a^j$ identifies $W^j$ with
$A^j$, and is an automorphism of $B$. The Fitting ideals are thus
constant in the stated range on a faithfully flat cover and hence
on the original parameter scheme.
\end{proof}

Combining this theorem with Lemma~\ref{lem:polarized-truncation}
gives the stable range $j\ge\min\{d-1,d-k+1\}$ for $d\ge2$.
The improvement is largest when the starting subspace has small
codimension. In particular, a codimension-two multiplication problem
is determined in degree three, however large $d$ is. Section~\ref{sec:full-codimension-two}
uses this ring-level assertion to retain the entire failure scheme,
including the closed locus where the defect is two.
